\documentclass[../main.tex]{subfiles}
\graphicspath{{\subfix{../images/}}}
\begin{document}

Die drei Gesamtlösungen aus dem morphologischen Kasten werden zunächst auf die Festforderungen geprüft. Wie bereits oben erwähnt schirmt Gesamtlösung 3 (Aluminium) das Bluetooth-Signal des ESP32 ab und verletzt damit Festforderung 5 (Funkdurchlässigkeit). Sie scheidet ohne weitere Bewertung aus. Die Gesamtlösungen 1 und 2 werden in einer Nutzwertanalyse (Tabelle \ref{tab:nwa}) mit einer Ideallösung verglichen. Dazu werden die Wunschforderungen mit den Gewichtungen aus Kapitel \ref{Rangfolgeverfahren} bewertet. Vergeben werden Punkte von 1 bis 10, die mit der jeweiligen Gewichtung multipliziert und aufsummiert werden. Die Ideallösung erreicht in jedem Kriterium 10 Punkte und damit einen gewichteten Gesamtwert von 10.

% ---------- Farben ----------
\definecolor{kopf}{RGB}{214,224,238}
\definecolor{ideal}{RGB}{226,239,218}
\definecolor{linie}{RGB}{120,120,120}

% ---------- Formatierung und Summen ----------
\newcommand{\fmt}[1]{\num{\fpeval{round(#1,3)}}}
\def\sumPa{0}\def\sumPb{0}\def\sumPi{0}
\def\sumGa{0}\def\sumGb{0}\def\sumGi{0}

% \zeile{Anforderung}{Gewichtung}{P Lösung 1}{P Lösung 2}
\newcommand{\zeile}[4]{%
  \xdef\sumPa{\fpeval{\sumPa+#3}}\xdef\sumPb{\fpeval{\sumPb+#4}}%
  \xdef\sumPi{\fpeval{\sumPi+10}}%
  \xdef\sumGa{\fpeval{\sumGa+#2*#3}}\xdef\sumGb{\fpeval{\sumGb+#2*#4}}%
  \xdef\sumGi{\fpeval{\sumGi+#2*10}}%
  #1 & \fmt{#2}
  & #3 & \fmt{#2*#3}
  & #4 & \fmt{#2*#4}
  & \cellcolor{ideal}10 & \cellcolor{ideal}\fmt{#2*10} \\ \hline}

% ---------- Gedrehte Tabelle, mittig auf der Seite ----------
\begin{center}
\begin{turn}{90}
\begin{minipage}{0.65\textheight}
\centering
\captionof{table}{Quantitative Nutzwertanalyse der Gesamtlösungen 1 und 2 im Vergleich zur Ideallösung}
\label{tab:nwa}
\vspace{6pt}
\adjustbox{max width=\linewidth}{%
\small\sffamily
\setlength{\tabcolsep}{3pt}
\arrayrulecolor{linie}
\renewcommand{\arraystretch}{1.7}
\begin{tabular}{|>{\centering\arraybackslash}p{3.4cm}|>{\centering\arraybackslash}p{2.2cm}|*{6}{>{\centering\arraybackslash}p{1.75cm}|}}
\hline
\rowcolor{white}
\multicolumn{2}{|c|}{} &
\multicolumn{2}{c|}{\cellcolor{kopf}\textbf{Gesamtlösung 1}} &
\multicolumn{2}{c|}{\cellcolor{kopf}\textbf{Gesamtlösung 2}} &
\multicolumn{2}{c|}{\cellcolor{kopf}\textbf{Ideallösung}} \\ \hline
\rowcolor{kopf}
\textbf{Anforderung} & \makecell{\textbf{Relative}\\\textbf{Gewichtung}} &
\makecell{\textbf{Punkte}\\\textbf{1--10}} & \makecell{\textbf{Punkte}\\\textbf{gewichtet}} &
\makecell{\textbf{Punkte}\\\textbf{1--10}} & \makecell{\textbf{Punkte}\\\textbf{gewichtet}} &
\makecell{\textbf{Punkte}\\\textbf{1--10}} & \makecell{\textbf{Punkte}\\\textbf{gewichtet}} \\ \hline

\zeile{Bauraum}                       {0.14}{10}{7}
\zeile{Lebensdauer}                   {0.15}{7}{8}
\zeile{DIY-Fertigungsverfahren}       {0.18}{10}{10}
\zeile{Montage ohne Spezialwerkzeug}  {0.13}{8}{8}
\zeile{Montage-/Demontagezeit}        {0.07}{9}{9}
\zeile{Gesamtgewicht}                 {0.04}{8}{6}
\zeile{Herstellungskosten}            {0.20}{8}{6}
\zeile{Fertigungstiefe}               {0.06}{3}{2}
\zeile{Recyclingfähigkeit}            {0.03}{9}{9}

\multicolumn{2}{|r|}{\textbf{Summe}\hspace{4pt}} &
\textbf{\fmt{\sumPa}} & \textbf{\fmt{\sumGa}} &
\textbf{\fmt{\sumPb}} & \textbf{\fmt{\sumGb}} &
\textbf{\fmt{\sumPi}} & \textbf{\fmt{\sumGi}} \\ \hline
\multicolumn{2}{|r|}{\textbf{Bezogene Werte}\hspace{4pt}} &
\textbf{\fmt{\sumPa/\sumPi}} & \textbf{\fmt{\sumGa/\sumGi}} &
\textbf{\fmt{\sumPb/\sumPi}} & \textbf{\fmt{\sumGb/\sumGi}} &
1 & 1 \\ \hline
\end{tabular}%
}
\end{minipage}
\end{turn}
\end{center}

Gesamtlösung 1 erreicht einen gewichteten Nutzwert von 8,29 und Gesamtlösung 2 von 7,48. Bezogen auf die Ideallösung entspricht das 0,829 bzw. 0,748. Auch ohne Gewichtung liegt Gesamtlösung 1 mit 72 von 90 Punkten (0,80) vor Gesamtlösung 2 mit 65 Punkten (0,72). Der Vorsprung entsteht vor allem beim Bauraum (10 gegenüber 7 Punkten), bei den Herstellungskosten (8 gegenüber 6) und beim Gesamtgewicht (8 gegenüber 6). Gesamtlösung 2 schneidet nur bei der Lebensdauer besser ab (8 gegenüber 7 Punkten). Dieser Vorteil fällt bei einer Gewichtung von 0,15 nicht ins Gewicht. Bei den übrigen Kriterien sind beide Lösungen gleichwertig bewertet.

Gesamtlösung 1 wird daher als Vorzugsvariante ausgewählt und in Kapitel \ref{Zeichnungssatz} als Zeichnungssatz ausgearbeitet.

\end{document}